% IDS Theory --- Consciousness Edition (formal paper format)
% Derived from IDS Canonical Theory v2.1.0-r2 (public-r6.2)
% Build: latexmk -pdf IDS_Consciousness_Edition_EN.tex
\documentclass[11pt,a4paper]{article}

\usepackage[T1]{fontenc}
\usepackage[utf8]{inputenc}
\usepackage{mathptmx}
\usepackage[scaled=0.92]{helvet}
\usepackage{microtype}
\usepackage{amsmath,amssymb,amsthm}
\usepackage{mathrsfs}
\usepackage[margin=2.6cm]{geometry}
\usepackage{booktabs,array,tabularx,xltabular}
\usepackage{ragged2e}
\usepackage{enumitem}
\usepackage{xcolor}
\usepackage[framemethod=TikZ]{mdframed}
\usepackage[colorlinks=true,linkcolor=linknavy,citecolor=linknavy,urlcolor=linknavy]{hyperref}

\definecolor{linknavy}{RGB}{31,79,111}
\definecolor{boxline}{RGB}{110,100,80}
\definecolor{boxbg}{RGB}{247,245,238}

\emergencystretch=1.5em
\setlength{\parskip}{0pt}

\newcolumntype{Y}{>{\RaggedRight\arraybackslash}X}

\mdfdefinestyle{canon}{%
  linecolor=boxline,linewidth=0.9pt,
  backgroundcolor=boxbg,
  topline=false,bottomline=false,rightline=false,
  innerleftmargin=10pt,innerrightmargin=10pt,
  innertopmargin=7pt,innerbottommargin=7pt,
  skipabove=10pt,skipbelow=10pt}
\newenvironment{canonbox}{\begin{mdframed}[style=canon]\itshape}{\end{mdframed}}
\newenvironment{defbox}[1]{\begin{mdframed}[style=canon]\textbf{\textsc{#1}.}\ \itshape}{\end{mdframed}}

\newenvironment{namedthm}[1]{\par\addvspace{\medskipamount}\noindent\textbf{#1.}\ \itshape\ignorespaces}{\par\addvspace{\medskipamount}}

\theoremstyle{plain}

\title{\bfseries What Prediction Leaves Behind\\[6pt]
\large An Information-Driven Stress account of memory, the self, and consciousness}
\author{\normalsize A consciousness-oriented edition prepared from the IDS Canonical Theory, revision public-r6.2}
\date{July 16, 2026}

\begin{document}

\maketitle

\begin{abstract}
\noindent
The predictive brain minimizes error; but error minimization does not, by itself, classify what a mind
becomes. Two systems can lower the same objective and yet differ in what they remember, avoid, or fix
in place. This essay develops a consciousness-oriented reading of Information-Driven Stress (IDS)
theory around that gap. Its central object is not a scalar quantity but an equivalence class: the
non-integrable \emph{structural residual} $[\omega]\in C^1/\operatorname{im}d_0$ --- the part of a
mind's activity that no local adjustment can dissolve. Memory, individuality, trauma, and habit are
read as residual structure that survived simplification rather than error that was eliminated. The
edition makes one claim sharply. The free-energy view builds the subject into its primitives:
``prediction error'' presupposes a predictor. IDS instead treats the observer as an explicit variable
and the self as a \emph{regeneration closure} --- a rule that re-identifies itself from the inside ---
and proves that this rule cannot be certified from the observed return dynamics alone. Consciousness is
therefore approached not as a solved mechanism but as an explicit construction problem with a stated
boundary. Nothing here claims that the physical implementation of consciousness has been demonstrated.
\end{abstract}

\medskip
\tableofcontents

\section{What does prediction leave behind?}\label{sec:question}

A powerful picture organizes much of contemporary neuroscience: the brain holds a generative model of
the world, predicts its sensory input, and revises perception and action to reduce the resulting
prediction error. Sharpened, this becomes the free-energy principle (FEP), under which perception,
learning, and action all descend a single variational functional. The picture explains a great deal.

Yet a clinician watching a brain recover from injury sees something the objective alone does not name.
After damage, a brain does not simply become a lower-performing version of itself. Some functions are
lost; some migrate; circuits that were quiet come forward; and sometimes repetitions, fixations,
sensations, or affects that did not exist in health settle in as a \emph{new stable state}. The brain
looks less like a machine erasing its errors and more like one \emph{reworking} them into different
structure.

This motivates a question that a scalar objective cannot answer on its own:

\begin{canonbox}
When a mind lowers its prediction error, what was actually resolved --- and what was merely moved,
hidden, externalized, or fixed in place as new structure?
\end{canonbox}

IDS theory is built to make that question precise. Its answer begins by refusing to identify the
content of a mind with the \emph{size} of its error.

\section{Two kinds of stress in a mind}\label{sec:stress}

IDS distinguishes two fundamentally different notions of stress. The first is a scalar: a potential
$\mathcal S$ on the space of a system's observable states, dissipated by gradient descent --- the
formal cousin of free energy. It answers ``how much tension is there?'' The second is not a number at
all.

Represent a mind's interactions on a finite oriented complex $X$, with a cochain sequence
$C^0(X)\xrightarrow{d_0}C^1(X)\xrightarrow{d_1}C^2(X)$ and $d_1d_0=0$; edges in $C^1$ carry flows of
belief-update, prediction-error propagation, and response. Any observed flow $\omega\in C^1$ splits by
the finite Hodge decomposition,
\begin{equation}
\omega=d_0\phi+h+\delta_1\psi,
\end{equation}
into a gradient part $d_0\phi$ (what a local potential can produce, hence dissipate), a harmonic part
$h$ (closed but not exact --- circulation with no potential), and a coexact part $\delta_1\psi$. What no
potential can generate is the \textbf{structural residual}, and its invariant is the quotient class
\begin{equation}
\boxed{\,[\omega]\in R^1(X):=C^1(X)/\operatorname{im}d_0\,.}
\label{eq:residual}
\end{equation}
It answers the second question: ``what shape does the irreducible tension have?'' A \emph{bug}, in this
language, is exactly a nonzero residual class relative to a given description --- not an absolute defect,
since a change of description may dissolve it or amplify it.

The consequence for the mind is immediate. \emph{Debugging} is not driving the residual to zero. It is
any reorganization of residual structure, and stabilizing a nonzero structure $\rho^\ast\ne0$ counts as
much as removing one. Recovery, compensation, adaptation, chronification, learning, and personality
change all become moves in the same space --- transformations of $[\omega]$ --- rather than degrees of a
single success measure.

\begin{canonbox}
$\mathcal S$ governs dissipation; $[\omega]$ governs structure. What a mind \emph{is} lives in the
residual, not in the magnitude of its error.
\end{canonbox}

\section{The subject is not a premise}\label{sec:subject}

Here the consciousness-oriented reading departs most sharply from the predictive-brain tradition. The
basic quantity of FEP is the prediction error $\varepsilon = o - g(\mu)$, the discrepancy between input
$o$ and what a generative map $g$ expects given internal state $\mu$. This quantity does not exist
without $\mu$: an expecting subject is required for ``error'' to be defined at all. The boundary that
separates inner from outer --- the Markov blanket --- is likewise assumed rather than derived. In short,
the free-energy view weaves the core of the subject into its primitives.

The residual class \eqref{eq:residual} is different. To say $[\omega]\ne0$ is to say that no potential
$\phi$ makes $\omega=d_0\phi$: a cohomological fact fixed by $d_0$ alone, requiring no expectation and
no subject. The observer enters only afterward, and as a \emph{variable}. An observer is a distribution
$\mu\in\mathcal P(\mathcal O)$ whose kernels and metric,
\begin{equation}
\langle K\rangle_\mu=\int K_o\,d\mu(o),\qquad g_\mu,
\end{equation}
fix the inner product. This yields a clean separation:
\begin{equation}
\underbrace{[\omega]\in C^1/\operatorname{im}d_0}_{\text{observer-independent}}
\qquad\text{versus}\qquad
\underbrace{\mathfrak r=h+\delta_1\psi}_{\text{observer-dependent reading}}.
\end{equation}
The \emph{existence} of the residual is objective, independent of any metric; only its \emph{reading}
--- how much looks like protected memory $h$ and how much like forced flux $\delta_1\psi$ --- is relative
to the observer. What the predictive view hides inside the word ``error,'' IDS extracts as an explicit
degree of freedom.

\section{The self as a regeneration closure}\label{sec:self}

If the subject is not a premise, what is it? IDS constructs the self not as a primitive but as a
\textbf{regeneration closure}. Call a system self-regenerating when, with external template input set to
zero, its autonomous update regenerates from its own internal state and resources not merely its state
but the machinery of self-identification itself: the realization and read-out of its residual
$(\iota,\varpi)$, the relevance projection $P_{\mathrm{ctrl}}$, the comparison metric $G$, the intrinsic
return event $\tau$, and the rule $\mathscr R$ by which it recognizes a later state as \emph{itself}. The
self, on this account, is the re-identification rule that survives the wholesale turnover of its own
parts.

Molecules, memories, values, and circuits all change; a self persists across that change not by
conserving a substance but by continuously rebuilding the rule that decides what counts as the same
self --- which changes to repair, which to accept, which to externalize, and where to draw its boundary.

IDS then imposes a discipline that the predictive view does not. Writing $\mathscr R=DF|_{x_\ast}$ for
the linearized return operator at a return fixed point, the following holds.

\begin{namedthm}{No-go theorem (endogenous vs.\ external self-constitution)}
An internal loop that regenerates a map $F$ and an external oracle that writes the same $F$ each cycle
produce identical observed return dynamics. Hence no functional of the return operator $\mathscr R$
alone can distinguish endogenous self-constitution from external reconstitution; the distinction
requires causal-boundary information beyond the operator.
\end{namedthm}

\begin{proof}[Sketch]
Both realizations share one $C^1$ input--output map $F$, so every invariant built from $F$ and its
derivatives coincides. Only an intervention that zeroes the external input separates them; therefore the
separation lies outside operator theory.
\end{proof}

The predictive view acquires the subject for free, by assumption. IDS records that the subject
\emph{cannot} be obtained for free from the observed dynamics, and specifies instead what must be added
--- externally template-free causal closure --- before a system's self-maintenance may be called its own.
For a theory that reaches toward consciousness, this is the honest posture: the self is placed at the
exit, as a conditional construction, not at the entrance as a given.

\section{What survives simplification}\label{sec:coarse}

No mind retains everything; it must coarse-grain. IDS does not treat coarse-graining as mere information
loss but as a rule for which differences may be identified. A coarse-graining $\Phi$ is
\emph{admissible} when it sends gradients to gradients, $\Phi(\operatorname{im}d_0)\subseteq
\operatorname{im}d_0$, and then it descends to the residual quotient and transports $[\omega]$ faithfully.
A coarse-graining that integrates out hidden degrees of freedom generally breaks this, and the defect
\begin{equation}
\mathfrak A_\Phi=q\circ\Phi\circ d_0\ne 0
\end{equation}
is an \emph{anomaly}: something microscopically exact reappears, after simplification, as genuine
residual. The forgetting is not free; its cost returns as structure.

The permissible boundary of forgetting has a crisp form: one may safely identify only those differences
that do not change future response under admissible intervention. This is the formal core of a familiar
clinical intuition --- that discarding a detail is harmless exactly when it never alters what the system
will later do or need.

Two consequences bear directly on mind. First, the arrow of time is a consequence of forgetting: perfect
memory is reversible and yields no macroscopic direction; it is because coarse-graining erases that an
entropy gradient, and a felt direction of time, exist. Second, the past is not rewritten but reread. The
residual class of a memory changes when the coarse-graining or the observation metric changes,
$[\omega_{\mathrm{past}}]_{X,g}\to[\omega_{\mathrm{past}}]_{X',g'}$; psychological integration and
reinterpretation have exactly this form. Forgetting, in IDS, is not a defect but a condition for the
generation of structure.

\section{Memory with a history}\label{sec:types}

Residual structure that persists comes in distinct kinds, and the distinctions map onto recognizable
features of mind (Table~\ref{tab:types}). Writing the equilibrium of the residual dynamics
$\dot\rho=-L\rho+a$ (with curvature Laplacian $L=\delta_1d_1$) as $\rho^\ast=h+L^{+}a$:

\begin{table}[htbp]
\centering\small
\begin{tabularx}{\textwidth}{@{}l Y Y@{}}
\toprule
Kind & Structure & Reading in a mind \\
\midrule
Type I & Topological: $\rho^\ast=h\in H^1$ & A memory or identity protected because it is a closed loop nothing local can cut \\
Type II & Forced: $\rho^\ast=L^{+}a$ & A state held only while input flows --- much of living, adaptive order \\
Type III & Promoted: $\mathcal T_{\mathrm{III}}=\bar H^1\ominus H^1$ & A once-transient constraint that has become a permanent trait \\
\bottomrule
\end{tabularx}
\caption{Three kinds of residual structure and their mental readings.}
\label{tab:types}
\end{table}

Type III is the subtle one. It is not a third component at a fixed scale but a \emph{genealogical} type:
a mode that decayed at the fine scale becomes harmonic --- protected --- at the coarse scale. When only
the finest constraints are forgotten, the promoted sector is exactly the orthogonal complement
$\bar H^1\ominus H^1$, and its canonical slogan reads
\begin{equation}
\boxed{\ \text{local fine curvature / constraint response}\ \longrightarrow\ \text{coarse global harmonic memory}.\ }
\end{equation}
A childhood constraint that is no longer active can nonetheless be locked into the adult as protected
structure. Trauma fits the Type I picture --- a closed-but-not-exact loop of memory that pure dissipation
cannot open; repetition is re-excitation of the same residual class; integration is re-representation of
that class under a new coarse-graining, a new language for the same history.

\section{Path-dependence: holonomy}\label{sec:holonomy}

Returning to the same observable state does not return a mind to the same internal state. Passing through
states and back, a system can respond differently than before: the path is written into the present rule,
not just the present value. IDS models this with holonomy. On a bundle of belief/scale fibers carrying a
comparison connection $A$, the transformation left after traversing a closed loop $\gamma$ is
\begin{equation}
\operatorname{Hol}_\gamma(A)=\mathcal P\exp\!\oint_\gamma A\ \ne\ \mathrm{id}.
\end{equation}
The base point returns; the fiber does not. Learned reactivity, post-traumatic vigilance, altered
sensitivity after exposure, motor control after repeated training, and the history-dependence of trust
and attachment are all describable this way: they cannot be read off the current stimulus, because the
\emph{order} in which a mind engaged the world is retained in how it now transforms. Holonomy is not
itself memory or awareness; it is the geometric vessel in which a history the present observation cannot
recover may nonetheless be held.

\section{The ledger of forgetting: information enthalpy}\label{sec:enthalpy}

When a mind simplifies, information does not vanish for free. The predictive view works throughout with
scalars --- free energy, surprisal, entropy --- but the existence of a global scalar potential is itself
a strong assumption. IDS proposes that the object governing hidden information exchange is not a scalar
but a connection: an exchange one-form $\Xi_I=d\mathcal H_I+W_I$, with information enthalpy
$\mathcal H_I$, whose kernel defines an information-adiabatic comparison. A scalar ``information
entropy'' is then well defined only under two conditions: flatness of this connection (the
integrability that lets an entropy state-function exist, the information-geometric analogue of the second
law's integrating factor), and calibration against a maximum-entropy reference. Even when flat, a global
holonomy can remain --- a \emph{thermal memory} that no local reading detects.

The reading for mind is that the entropy of a subjective state is a derived, conditional quantity, not a
primitive; and that the cost of what a mind discards is generally carried not by a number but by the
curvature and holonomy of this exchange. Uncertainty, bias, symptom, and bodily response are where the
discarded degrees of freedom return.

\section{Resilience of the self}\label{sec:resilience}

Near a stable structure, linearized recovery obeys $\|y(t)\|\le e^{-\lambda^{+}_{\min}(L)\,t}\|y(0)\|$,
so the positive spectral gap $\Delta_{\mathrm{res}}=\lambda^{+}_{\min}(L)$ measures how fast residual
structure returns after perturbation. But directions with no direct dissipation can still recover, if a
conservative rotation keeps turning them into dissipative ones until the whole space is reached --- a
\emph{hypocoercive} resilience that formalizes the capacity to route damage that cannot be repaired
directly toward a channel where it can. Clinically, this is why a recovered score must be read with care:
the same improvement can be genuine repair of the original circuit, compensation by another, a changed
strategy, avoidance of the eliciting situation, or a deficit invisible to the measure. What must be
tracked is not whether the error fell but which structure formed while it fell.

\section{What this does and does not claim}\label{sec:claims}

This edition is an interpretive and mathematical framework, and its honesty is part of its content. It
does \emph{not} claim that the physical implementation of consciousness has been demonstrated, nor that
awareness has been reduced to residual cohomology. The clinical and psychological readings above are
interpretations, disciplined by the mathematics but not theorems. Holonomy is offered as a vessel for
history, not as an account of experience. The self is offered as a regeneration closure whose endogenous
character, by the no-go theorem, cannot be certified from observed dynamics alone.

The framework is nonetheless answerable. It earns an independent place only if systems that lower the
same objective can be told apart by their residual structure: if the future response of equally
optimized minds is predicted by their harmonic content, boundary flow, and holonomy; if locally resolved
inconsistencies re-emerge at another level, in the body, or in the environment; and if forgetting
information that changes future control produces long-run breakdown despite short-run fit. Where a
residual measure adds no prediction beyond error, surprisal, and entropy, the extra apparatus should be
discarded. Where it does, it names something the scalar view cannot.

\section{Conclusion}\label{sec:conclusion}

The predictive brain lowers the error it can lower. What it cannot lower it carries, and where it carries
its irreducible tension --- as protected loops, as promoted constraints, as path-written holonomy, as the
returning cost of what was forgotten --- is where memory, individuality, and symptom take their shape.

\begin{canonbox}
A mind is not a machine that erases its errors. It reduces the errors it can, and transports, remembers,
and re-forms the residues it cannot --- and the self is not the result of that process but the rule that
keeps re-identifying itself while the process runs.
\end{canonbox}

The heart of the account is a single object,
\[
\boxed{\,[\omega]\in C^1/\operatorname{im}d_0\,,}
\]
the residue that no local correction dissolves. What we are made of is not only the information a mind
could keep, but also what it could not erase --- reworked into new form, and carried forward as part of
the self. Whether that self is conscious, and how, is left here as an open and explicitly bounded
question, not a settled one.

\section*{A note on this edition}
\addcontentsline{toc}{section}{A note on this edition}

This consciousness-oriented edition is a condensed, general-audience derivative of the IDS Canonical
Theory v2.1.0-r2 (public revision r6.2). It re-centers the framework on mind, memory, and the self,
keeping the residual core, the observer as an explicit variable, the regeneration closure with its no-go
boundary, and the falsifiability commitments, while omitting the full proofs, audit ledgers, and the
physical-bridge and gauge-theoretic layers. The canonical document remains the authoritative source.

\bigskip
\noindent\footnotesize\textbf{Selected references.}
K.~Friston, \emph{The free-energy principle: a rough guide to the brain}, Trends Cogn.\ Sci.\ (2009);
S.~J.~Gershman, \emph{What does the free energy principle tell us about the brain?}, Neural Comput.\ (2019);
T.~Parr, L.~Da~Costa, K.~Friston, \emph{Markov blankets, information geometry and stochastic thermodynamics},
Phil.\ Trans.\ R.\ Soc.\ A (2020);
V.~Raja et al., \emph{The Markov blanket trick: on the scope of the free energy principle and active inference},
Phys.\ Life Rev.\ (2021);
M.~J.~D.~Ramstead et al., \emph{On Bayesian mechanics: a physics of and by beliefs}, Interface Focus (2023).

\end{document}
